% !TeX root = thesis.tex

\chapter{Introduction} \label{chap:intro}

\section*{}
 

\section{Context} \label{sec:context}

The 21st century has witnessed remarkable advances in Artificial Intelligence (AI) systems, reshaping numerous industrial and societal sectors. (\cite{Golec2025}). 
The introduction of AlexNet in 2012 (\cite{Alex2012}), AlphaGo's victory over Lee Sedol in 2016 (\cite{Silver2016}),
and the public release of ChatGPT in 2022 (\cite{OpenAI2024}) are among the most significant milestones in modern AI,
advancing the fields of Computer Vision (CV), Sequential Decision-Making, and Natural Language
Processing (NLP), respectively. These breakthroughs were driven by a rapid growth in
computational power, the development of specialized hardware such as GPUs and TPUs,
and the widespread adoption of Deep Learning (DL) techniques (\cite{Sajadieh2026}).

As a result, DL has become the dominant paradigm in modern Machine Learning (ML) due to its ability to
learn highly complex and abstract representations directly from data (\cite{Lecun2015}).

\section{AI in Industry 4.0} \label{sec:aiindustry}

Over the last fifteen years, the industrial sector has undergone a profound digital transformation,
leading to the emergence of the Industry 4.0 paradigm. Industry 4.0 is characterized by the integration
of advanced manufacturing technologies, Cyber-Physical Systems (CPS), the Internet of
Things (IoT), cloud computing, and interconnected production environments. These technologies
enable the continuous acquisition of large volumes of sensor data, creating smart factories capable of
monitoring, communicating, and adapting their operation in real time.

The availability of these large-scale industrial datasets has created favourable conditions for the
widespread adoption of DL techniques in this field. AI systems are increasingly being
employed to support a variety of tasks, such as production planning and scheduling, process monitoring and control, Predictive
Maintenance (PdM) and robotics, contributing to improved efficiency,
reliability, and productivity in modern manufacturing environments (\cite{Cerqueus2025}).

Among these applications, three closely related tasks are central to the monitoring of industrial 
equipment. Anomaly Detection (AD) is concerned with identifying observations that deviate from the 
expected behaviour of a system, without necessarily establishing what caused the deviation. Fault 
Detection (FD) goes one step further, determining whether a fault is present and, where possible, which 
type of fault is responsible for it. PdM extends the analysis in time, estimating the future evolution of 
the equipment's condition so that maintenance can be scheduled before a failure occurs. Together they 
describe successive stages of the same process, from noticing that something is wrong, to identifying the cause and 
to anticipate its consequences, and all three are addressed jointly in this dissertation.

\section{Motivation} \label{sec:motivation}

Biomethane upgrading plants operate continuously, and their output is typically committed to a gas grid 
or to a downstream customer under contract. In the Vacuum Pressure Swing Adsorption (VPSA) process 
considered in this dissertation, the vacuum pump is a single point of failure: it is what draws the vacuum that regenerates the adsorption beds between cycles, 
so if it stops, the entire plant stops with it. An unplanned failure 
therefore does not degrade production, it interrupts it completly, and the resulting downtime lasts as long as it 
takes to diagnose the problem, obtain the parts and carry out the repair.

Maintenance of this equipment is currently driven either by the calendar or by the failure itself: 
components are replaced at the fixed intervals recommended by the manufacturer, regardless of their actual 
condition, or repaired once they have already broken down. The motivation for this work is to move that 
decision towards the condition of the equipment itself. If degradation can be recognised from the data the 
plant already collects, while it is still developing, maintenance can be planned for a moment that suits 
production rather than forced by a breakdown, and interventions can be limited to the components that 
genuinely require them. The intended outcome is a plant that runs with fewer unplanned interruptions, 
lower maintenance costs, and a reduced likelihood of a minor deviation developing into a critical event.

\section{Problem Definition} \label{sec:problem}

Three problems stand in the way of this condition-based approach.

The first is that the plant already has the sensors, but only uses them for fixed alarm limits. The
manufacturer sets a warning and a trip value for each variable, and the control system reacts when one of
them is crossed. By that point the fault is usually well developed. The alarm also tells the operator only
that a limit was passed, not which fault caused it or how far it has progressed, and it looks at one
sensor at a time, while a real fault affects several at once.

The second is that the behaviour of the pump is too complex to be captured by fixed rules. Its fourteen
monitored variables are not independent, so a developing fault rarely shows up as a single variable
crossing a line. Instead, it shows up as a change in how those variables move in relation to one another,
and in how that pattern evolves over the following hours. Expressing this as a set of rules would require
knowing every relevant interaction and every fault signature in advance, and the resulting rules would be
numerous and fragile. DL offers a way around this, since it learns these relationships, and the
way they change over time directly from the data rather than requiring them to be specified beforehand.

The third is that there is no fault data. The company has never recorded a vacuum pump failure at any of
its plants. That is good news operationally, but it leaves nothing to learn from: a supervised model needs
examples of what it is meant to detect, and without them there is also no way to measure how well it
works. The data has to be created before anything can be trained on it.

This dissertation addresses these problems by building a labelled dataset through controlled fault
injection, and then developing DL models that perform AD, FD and
PdM together, using the sensor data the plant already collects.

\section{Objectives} \label{sec:objectives}

The objectives of this thesis are:

\begin{itemize}
    \item To perform an Exploratory Data Analysis (EDA) of the biomethane upgrading dataset in order to better understand the data, characterize the process variables and identify potential challenges for model development.
    \item To develop and train single-task DL models for AD, FD, and PdM.
    \item To develop and train multi-task DL models using multiple architectural configurations capable of jointly performing several prediction tasks.
    \item To develop and train traditional ML models for each prediction task, providing a baseline for comparison.
    \item To compare and evaluate the performance of all proposed models using appropriate evaluation metrics.
    \item To investigate whether multi-task learning improves predictive performance compared with single-task DL models and traditional ML approaches.
\end{itemize}

\section{Company Presentation} \label{sec:company}

This dissertation was developed in collaboration with Sysadvance, a Portuguese company specialized in gas separation and purification technologies. The company 
develops industrial solutions for gas generation, purification, and biogas upgrading, providing the real industrial case study considered throughout this work.

\section{Thesis Structure} \label{sec:thesis}

This dissertation is organized into 7 chapters.

\begin{itemize}

    \item \textbf{Chapter 1} introduces the research context, motivation, objectives, and overall organization of the dissertation.
    \item \textbf{Chapter 2} presents the theoretical background required to understand the problem addressed in this work. It introduces the concepts of biogas 
and biomethane production, biogas upgrading technologies, the considered Vacuum Pressure Swing Adsorption (VPSA) plant, the operating principle of the vacuum 
pump, and the monitored instrumentation.
    \item \textbf{Chapter 3} reviews the theoretical foundations and the state of the art, covering artificial
    intelligence, machine learning and deep learning, the industrial monitoring and maintenance strategies
    addressed in this work, and the related literature that defines the research gap.

    \item \textbf{Chapter 4} describes the methodology, including the dataset, the exploratory data analysis,
    the data cleaning pipeline, the fault injection framework, the compilation of the training and test sets,
    the task formulation, and the models developed.

    \item \textbf{Chapter 5} presents the experimental results obtained for the three model families across
    the three tasks and the window sizes tested.

    \item \textbf{Chapter 6} discusses those results, and states the limitations of the study.

    \item \textbf{Chapter 7} presents the conclusions and outlines directions for future work.

\end{itemize}